%!TeX root=../princesstop.tex
\chapter{The Other Side of the Wall}

\begin{figure}[t!]
\centering
\includegraphics[width=\linewidth]{filigree}
\end{figure}

\lettrine[lines=5]{W}{hen} one lives in a row of houses, it is interesting to think of the things which are being done and said on the other side of the wall of the very rooms one is living in. Sara was fond of amusing herself by trying to imagine the things hidden by the wall which divided the Select Seminary from the Indian gentleman's house. She knew that the schoolroom was next to the Indian gentleman's study, and she hoped that the wall was thick so that the noise made sometimes after lesson hours would not disturb him.

<I am growing quite fond of him,> she said to Ermengarde; <I should not like him to be disturbed. I have adopted him for a friend. You can do that with people you never speak to at all. You can just watch them, and think about them and be sorry for them, until they seem almost like relations. I'm quite anxious sometimes when I see the doctor call twice a day.>

<I have very few relations,> said Ermengarde, reflectively, <and I'm very glad of it. I don't like those I have. My two aunts are always saying, <Dear me, Ermengarde! You are very fat. You shouldn't eat sweets,> and my uncle is always asking me things like, <When did Edward the Third ascend the throne?> and, <Who died of a surfeit of lampreys?>>

Sara laughed.

<People you never speak to can't ask you questions like that,> she said; <and I'm sure the Indian gentleman wouldn't even if he was quite intimate with you. I am fond of him.>

She had become fond of the Large Family because they looked happy; but she had become fond of the Indian gentleman because he looked unhappy. He had evidently not fully recovered from some very severe illness. In the kitchen—where, of course, the servants, through some mysterious means, knew everything—there was much discussion of his case. He was not an Indian gentleman really, but an Englishman who had lived in India. He had met with great misfortunes which had for a time so imperilled his whole fortune that he had thought himself ruined and disgraced forever. The shock had been so great that he had almost died of brain fever; and ever since he had been shattered in health, though his fortunes had changed and all his possessions had been restored to him. His trouble and peril had been connected with mines.

<And mines with diamonds in 'em!> said the cook. <No savin's of mine never goes into no mines—particular diamond ones>—with a side glance at Sara. <We all know somethin' of \textsc{them}.>

<He felt as my papa felt,> Sara thought. <He was ill as my papa was; but he did not die.>

So her heart was more drawn to him than before. When she was sent out at night she used sometimes to feel quite glad, because there was always a chance that the curtains of the house next door might not yet be closed and she could look into the warm room and see her adopted friend. When no one was about she used sometimes to stop, and, holding to the iron railings, wish him good night as if he could hear her.

<Perhaps you can \textsc{feel} if you can't hear,> was her fancy. <Perhaps kind thoughts reach people somehow, even through windows and doors and walls. Perhaps you feel a little warm and comforted, and don't know why, when I am standing here in the cold and hoping you will get well and happy again. I am so sorry for you,> she would whisper in an intense little voice. <I wish you had a <Little Missus> who could pet you as I used to pet papa when he had a headache. I should like to be your <Little Missus> myself, poor dear! Good night—good night. God bless you!>

She would go away, feeling quite comforted and a little warmer herself. Her sympathy was so strong that it seemed as if it MUST reach him somehow as he sat alone in his armchair by the fire, nearly always in a great dressing gown, and nearly always with his forehead resting in his hand as he gazed hopelessly into the fire. He looked to Sara like a man who had a trouble on his mind still, not merely like one whose troubles lay all in the past.

<He always seems as if he were thinking of something that hurts him \textsc{now},> she said to herself, <but he has got his money back and he will get over his brain fever in time, so he ought not to look like that. I wonder if there is something else.>

If there was something else—something even servants did not hear of—she could not help believing that the father of the Large Family knew it—the gentleman she called Mr Montmorency. Mr Montmorency went to see him often, and Mrs Montmorency and all the little Montmorencys went, too, though less often. He seemed particularly fond of the two elder little girls—the Janet and Nora who had been so alarmed when their small brother Donald had given Sara his sixpence. He had, in fact, a very tender place in his heart for all children, and particularly for little girls. Janet and Nora were as fond of him as he was of them, and looked forward with the greatest pleasure to the afternoons when they were allowed to cross the square and make their well-behaved little visits to him. They were extremely decorous little visits because he was an invalid.

<He is a poor thing,> said Janet, <and he says we cheer him up. We try to cheer him up very quietly.>

Janet was the head of the family, and kept the rest of it in order. It was she who decided when it was discreet to ask the Indian gentleman to tell stories about India, and it was she who saw when he was tired and it was the time to steal quietly away and tell Ram Dass to go to him. They were very fond of Ram Dass. He could have told any number of stories if he had been able to speak anything but Hindustani. The Indian gentleman's real name was Mr Carrisford, and Janet told Mr Carrisford about the encounter with the little-girl-who-was-not-a-beggar. He was very much interested, and all the more so when he heard from Ram Dass of the adventure of the monkey on the roof. Ram Dass made for him a very clear picture of the attic and its desolateness—of the bare floor and broken plaster, the rusty, empty grate, and the hard, narrow bed.

<Carmichael,> he said to the father of the Large Family, after he had heard this description, <I wonder how many of the attics in this square are like that one, and how many wretched little servant girls sleep on such beds, while I toss on my down pillows, loaded and harassed by wealth that is, most of it—not mine.>

<My dear fellow,> Mr Carmichael answered cheerily, <the sooner you cease tormenting yourself the better it will be for you. If you possessed all the wealth of all the Indies, you could not set right all the discomforts in the world, and if you began to refurnish all the attics in this square, there would still remain all the attics in all the other squares and streets to put in order. And there you are!>

Mr Carrisford sat and bit his nails as he looked into the glowing bed of coals in the grate.

<Do you suppose,> he said slowly, after a pause—<do you think it is possible that the other child—the child I never cease thinking of, I believe—could be—could POSSIBLY be reduced to any such condition as the poor little soul next door?>

Mr Carmichael looked at him uneasily. He knew that the worst thing the man could do for himself, for his reason and his health, was to begin to think in the particular way of this particular subject.

<If the child at Madame Pascal's school in Paris was the one you are in search of,> he answered soothingly, <she would seem to be in the hands of people who can afford to take care of her. They adopted her because she had been the favourite companion of their little daughter who died. They had no other children, and Madame Pascal said that they were extremely well-to-do Russians.>

<And the wretched woman actually did not know where they had taken her!> exclaimed Mr Carrisford.

Mr Carmichael shrugged his shoulders.

<She was a shrewd, worldly Frenchwoman, and was evidently only too glad to get the child so comfortably off her hands when the father's death left her totally unprovided for. Women of her type do not trouble themselves about the futures of children who might prove burdens. The adopted parents apparently disappeared and left no trace.>

<But you say <\textsc{if}> the child was the one I am in search of. You say <if.> We are not sure. There was a difference in the name.>

<Madame Pascal pronounced it as if it were Carew instead of Crewe—but that might be merely a matter of pronunciation. The circumstances were curiously similar. An English officer in India had placed his motherless little girl at the school. He had died suddenly after losing his fortune.> Mr Carmichael paused a moment, as if a new thought had occurred to him. <Are you \textsc{sure} the child was left at a school in Paris? Are you sure it was Paris?>

<My dear fellow,> broke forth Carrisford, with restless bitterness, <I am \textsc{sure} of nothing. I never saw either the child or her mother. Ralph Crewe and I loved each other as boys, but we had not met since our school days, until we met in India. I was absorbed in the magnificent promise of the mines. He became absorbed, too. The whole thing was so huge and glittering that we half lost our heads. When we met we scarcely spoke of anything else. I only knew that the child had been sent to school somewhere. I do not even remember, now, \textsc{how} I knew it.>

He was beginning to be excited. He always became excited when his still weakened brain was stirred by memories of the catastrophes of the past.

Mr Carmichael watched him anxiously. It was necessary to ask some questions, but they must be put quietly and with caution.

<But you had reason to think the school \textsc{was} in Paris?>

<Yes,> was the answer, <because her mother was a Frenchwoman, and I had heard that she wished her child to be educated in Paris. It seemed only likely that she would be there.>

<Yes,> Mr Carmichael said, <it seems more than probable.>

The Indian gentleman leaned forward and struck the table with a long, wasted hand.

<Carmichael,> he said, <I \textsc{must} find her. If she is alive, she is somewhere. If she is friendless and penniless, it is through my fault. How is a man to get back his nerve with a thing like that on his mind? This sudden change of luck at the mines has made realities of all our most fantastic dreams, and poor Crewe's child may be begging in the street!>

<No, no,> said Carmichael. <Try to be calm. Console yourself with the fact that when she is found you have a fortune to hand over to her.>

<Why was I not man enough to stand my ground when things looked black?> Carrisford groaned in petulant misery. <I believe I should have stood my ground if I had not been responsible for other people's money as well as my own. Poor Crewe had put into the scheme every penny that he owned. He trusted me—he \textsc{loved} me. And he died thinking I had ruined him—I—Tom Carrisford, who played cricket at Eton with him. What a villain he must have thought me!>

<Don't reproach yourself so bitterly.>

<I don't reproach myself because the speculation threatened to fail—I reproach myself for losing my courage. I ran away like a swindler and a thief, because I could not face my best friend and tell him I had ruined him and his child.>

The good-hearted father of the Large Family put his hand on his shoulder comfortingly.

<You ran away because your brain had given way under the strain of mental torture,> he said. <You were half delirious already. If you had not been you would have stayed and fought it out. You were in a hospital, strapped down in bed, raving with brain fever, two days after you left the place. Remember that.>

Carrisford dropped his forehead in his hands.

<Good God! Yes,> he said. <I was driven mad with dread and horror. I had not slept for weeks. The night I staggered out of my house all the air seemed full of hideous things mocking and mouthing at me.>

<That is explanation enough in itself,> said Mr Carmichael. <How could a man on the verge of brain fever judge sanely!>

Carrisford shook his drooping head.

<And when I returned to consciousness poor Crewe was dead—and buried. And I seemed to remember nothing. I did not remember the child for months and months. Even when I began to recall her existence everything seemed in a sort of haze.>

He stopped a moment and rubbed his forehead. <It sometimes seems so now when I try to remember. Surely I must sometime have heard Crewe speak of the school she was sent to. Don't you think so?>

<He might not have spoken of it definitely. You never seem even to have heard her real name.>

<He used to call her by an odd pet name he had invented. He called her his <Little Missus.> But the wretched mines drove everything else out of our heads. We talked of nothing else. If he spoke of the school, I forgot—I forgot. And now I shall never remember.>

<Come, come,> said Carmichael. <We shall find her yet. We will continue to search for Madame Pascal's good-natured Russians. She seemed to have a vague idea that they lived in Moscow. We will take that as a clue. I will go to Moscow.>

<If I were able to travel, I would go with you,> said Carrisford; <but I can only sit here wrapped in furs and stare at the fire. And when I look into it I seem to see Crewe's gay young face gazing back at me. He looks as if he were asking me a question. Sometimes I dream of him at night, and he always stands before me and asks the same question in words. Can you guess what he says, Carmichael?>

Mr Carmichael answered him in a rather low voice.

<Not exactly,> he said.

<He always says, <Tom, old man—Tom—where is the Little Missus?>< He caught at Carmichael's hand and clung to it. >I must be able to answer him—I must!< he said. >Help me to find her. Help me.>


On the other side of the wall Sara was sitting in her garret talking to Melchisedec, who had come out for his evening meal.

<It has been hard to be a princess today, Melchisedec,> she said. <It has been harder than usual. It gets harder as the weather grows colder and the streets get more sloppy. When Lavinia laughed at my muddy skirt as I passed her in the hall, I thought of something to say all in a flash—and I only just stopped myself in time. You can't sneer back at people like that—if you are a princess. But you have to bite your tongue to hold yourself in. I bit mine. It was a cold afternoon, Melchisedec. And it's a cold night.>

Quite suddenly she put her black head down in her arms, as she often did when she was alone.

<Oh, papa,> she whispered, <what a long time it seems since I was your <Little Missus>!>

This was what happened that day on both sides of the wall.